\documentclass[10pt]{article}

\usepackage[utf8]{inputenc}

\usepackage[T1]{fontenc}

\usepackage{amsmath}

\usepackage{amsfonts}

\usepackage{amssymb}

\usepackage[version=4]{mhchem}

\usepackage{stmaryrd}

\usepackage{bbold}



\begin{document}

многообразии $G$, типичными реализациями которой являются обобщенные функиии, заданные на этом многообразии $G$. Точнее, пусть $G$ - бесконечно гладкое многообразие и $D(G)$ - пространство бесконечно дифференцируемых финитных функций, определенных на $G$, с обычной топологией равномерной сходимости последовательностей равномерно финитных функций и всех их производных. Тогда на $G$ определено О.с. п., если задано непрерывное линейное отображение



\[

D(G) \rightarrow L_{0}(\Omega, \mathfrak{B}, \mu), \varphi \rightarrow f_{\varphi}, \varphi \in D(G),

\]



пространства $D(G)$ в пространство $L_{0}(\Omega, \mathfrak{B}, \mu)$ случайных величин, определенных на нек-ром вероятностном пространстве $\Omega$ с выделенной $\sigma$-алгеброй его подмножеств $\mathfrak{B}$ и вероятностной мерой $\mu$, определенной на $\mathfrak{B} ; L_{0}(\Omega, \mathfrak{B}, \mu)$ снабжено топологией сходимости по мере (см. [7]). В случае, когда вероятностным пространством является пространство $D^{\prime}(G)$ обобщенных функций на $G$ с $\sigma$-алгеброй $\mathfrak{B}_{0}$, порожденной цилиндрич. множествами в $D^{\prime}(G)$, а отображение задается формулами:



\[

f_{\varphi}(T)=(T, \varphi), T \in D^{\prime}(G), \varphi \in D(G),

\]



О. с. п. $\left\{f_{\varphi}, \varphi \in D^{\prime}(G)\right\}$ называется ка н о н и ческ им. Оказывается, что любое О. с. П. на конечномерном многообразии $G$ вероятностно изоморфно нек-рому (единственному) канонич. случайному полю на $G$ (см. [2]).



Приведенное здесь определение допускает ряд естественных модификаций: напр., можно рассмотреть О. с. п. с векторными значениями или вместо пространства $D(G)$ использовать в определении какое-нибудь более обширное пространство основных функций на $G$ (скажем, в случае $G=\mathbb{R}^{n}, n=1,2,3$, пространства $S\left(\mathbb{R}^{n}\right)$ бесконечно дифференцируемых функций, убывающих на бесконечности вместе со всеми производными, быстрее любой отрицательной степени $\left.|x|^{k}, k=-1,-2,-3, \ldots, x \in \mathbb{R}^{n}\right)$.



Понятие О. с. П. включает в себя классич. случайные поля и случайные процессы, реализациями к-рых являются обычные функции. Это понятие возникло в середине 50-х гг., когда обнаружилось, что многие естественные стохастич. образования не могут быть достаточно просто выражены в терминах классич. случайных полей, а на языке О.с. п. имеют простое и изящное описание. Так, напр., любая положительно определенная билинейная форма на $D\left(\mathbb{R}^{n}\right), n=1,2, \ldots$ :



\[

\left(\varphi_{1}, \varphi_{2}\right)=\int_{\mathbb{R}^{n}} \int_{\mathbb{R}^{n}} W\left(x_{1}, x_{2}\right) \varphi\left(x_{1}\right) \varphi\left(x_{2}\right) d x_{1} d x_{2},

\]



где $\varphi_{1}, \varphi_{2} \in D\left(\mathbb{R}^{n}\right), W\left(x_{1}, x_{2}\right)$ - положительно определенная симметрич. обобщенная функция двух переменных, однозначно определяет гауссовское О.с. п. $\left\{f_{\varphi}, \varphi \in D\left(\mathbb{R}^{n}\right)\right\}$ на $\mathbb{R}^{n}$ (с нулевым средним) так, что ковариация этого поля



\[

\int f_{\varphi_{1}} f_{\varphi_{2}} d \mu=\left(\varphi_{1}, \varphi_{2}\right)

\]



( $\mu$ - соответствующая этому полю вероятностная мера в $\left.D^{\prime}\left(\mathbb{R}^{n}\right)\right)$. Это О.с. п. оказывается классическим лишь при достаточно хорошей функции $W\left(x_{1}, x_{2}\right)$ (напр., непрерывной и ограниченной). Другие примеры: О.с.п. на $\mathbb{R}^{n}$ с независимыми значениями (см. [2]) или так наз. автомодельные случайные поля на $\mathbb{R}^{n}$ (см. [6]), среди к-рых вообще нет классич. полей.



Интерес к исследованию О.с. П. (особенно марковских случайных полей) возрос в последнее время из-за обнаруженной в начале 70-х гг. связи между задачей построения квантового физич. поля и построением марковских О.с. п. на $\mathbb{R}^{n}$ при $n>1$ (см. [5]).



Лит.: [1] Гельфанд И. М., Ш илов Г. Е., Пространства основных и обобщенных функций, М., 1958; [2] Гельфанд И. М., В и л е н к и н Н. Я., Некоторые применения гармонического анализа. Оснащенные гильбертовы пространства, М., 1961; [3] Гельф а нд И. М., «Доклады АН СССР», 1955, т. 100, № 5, с. 853-56;



398 ОБОБЩЕННОЕ [4] It o Kiyosi, «Mem. Coll. Sci. Univ. Kyoto. Ser. A», 1954, v. 28, № 3, p. 209-23; [5] С а й м о н Б., Модель $P(\varphi)_{2}$ евклидовой квантовой теории поля, пер. с англ., М., 1976; [6] Д о б р уш и н Р. Л., в сб.: Многокомпонентные случайные системы, М.- Л., 1978, с. 179-213; [7] Д о бр уш и н Р. Л., М и н л о с Р. А., «Успехи математич. наук», 1977 , т. 32, в. 2, с. 67-122. P. A. Минлос. ОБОБЩЕННЫЕ ИМПУЛЬСЫ - физические величинЫ $p_{i}$, определяемые формулами:



\[

p_{i}=\partial T / \partial \dot{q}_{i} \text { или } p_{i}=\partial L / \partial \dot{q}_{i} \text {, }

\]



где $T-$ кинетическая энергия, а $L-$ Лагранжа функция данной механической системы, выраженные через обобщенные координаты $q_{i}$ и обобщенные скорости $\dot{q}_{i}$. Размерность О. и. зависит от размерности обобщенной координаты. Если $q_{i}$ имеет размерность длины, то $p_{i}-$ размерность обычного импульса, то есть произведения массы на скорость; если же координатой $q_{i}$ является угол (величина безразмерная), то $p_{i}$ имеет размерность момента количества движения, и т. д.



ОБОБЩЕННЫЕ КООРДИНАТЫ - независимЫе параметры $q_{i}, i=1,2, \ldots, s$, любой размерности, число которых равно числу $s$ степеней свободы механической системы и которые однозначно определяют положение системы. Закон движения системы в О. к. дается $s$ уравнениями вида



\[

q_{i}=q_{i}(t)

\]



где $t$ - время. О. к. пользуются при решении многих задач, особенно когда система подчинена связям, налагающим ограничения на ее движение. При этом значительно уменьшается число уравнений, описывающих движение системы, по сравнению, напр., с уравнениями в декартовых координатах. В системах с бесконечно большим числом степеней свободы (сплошные среды, физич. поля) О.к. являются особые функции пространственных координат и времени, называемые потенциалами, волновыми функциями и т. П.



ОБОБЩЕННЫЕ ПОЛЯРНЫЕ КООРДИНАТЫ - СМ. Кoординаты.



ОБОБЩЕННЫЕ СиЛЫ - величины, играющие роль обычных сил, когда при изучении равновесия или движения механической системы ее положение определяется обобщенными координатами. Число О.с. равно числу $s$ степеней свободы системы; при этом каждой обобщенной координате $q_{i}$ соответствует своя О.с. $Q_{i}$. Значение О.с. $Q_{1}$, соответствующей координате $q_{1}$, можно найти, вычислив элементарную работу $\delta A_{1}$ всех сил на возможном перемещении системы, при к-ром изменяется только координата $q_{1}$, получая приращение $\delta q_{1}$. Тогда $\delta A_{1}=Q_{1} \delta q_{1}$, то есть коэффициент при $\delta q_{1}$ в выражении $\delta A_{1}$ и будет О. с. $Q_{1}$. Аналогично вычисляются $Q_{2}, Q_{3}, \ldots, Q_{s}$.



Размерность О.с. зависит от размерности обобщенной координаты. Если $q_{i}$ имеет размерность длины, то $Q_{i}-$ размерность обычной силы; если $q_{i}-$ угол, то $Q_{i}$ имеет размерность момента силы, и т. д. При изучении движения механич. системы О. с. входят вместо обычных сил в Лагранжа уравнения механики, а при равновесии все О. с. равны нулю. ОБОБЩЕННЫЙ ВЕКТОР СОСТОЯНИЯ - см. Вектор состояния.



ОБОБЩЕННЫЙ ИМПУЛЬС - ковектор на конфигурационном пространстве $X$ механической системы, то есть линейная функция на касательном пространстве к $X$. Если $q-$ координата на $X$, то $d q$ называется и м п у л ь с о м, соответствующим координате $q . \quad M$. Э. Казарян. ОБОБЩЕННЫЙ СЛУЧАЙНЫЙ ПРОЦЕСС - СМ. ОбОбщенное случайное поле, Случайный процесс.



ОБОБЩЕННЫЙ СОБСТВЕННЫЙ ВЕКТОР - $\mathrm{CM}$. Спектр оператора.



ОБОБЩЕННЫЙ XАРАКTEP п р е Д с т а в Л е н и Я - обобщенная функция $\chi(g)$ на $G$, определяемая равенством



\[

\int_{G} \chi(g) f(g) d g=\operatorname{tr} \int_{G} f(g) T_{g} d g .

\]





\end{document}